\subsection{\pbesym Terminates} \label{sec:pbe:termination}
\NV{HERE HERE HERE}
%% \subsection{Preliminaries} \label{sec:pbe:prelim}
%
Figure~\ref{fig:pbe:syntax} describes the syntax of
the main ingredients.
%
\RJ{describe $\Renv$}



\mypara{Normal}
%
We say $\Penv$ is \emph{normal}
if all applications in $\Penv$ are
of the form $\fapp{f}{x}$.
%
We can easily \emph{normalize}
a $\Penv$ by introducing fresh
variables for each sub-term, and
hence in the sequel we assume
that $\pbe{\Renv}{\Penv}{\pred}$
is only invoked with normal $\Penv$.

\mypara{Stores and Inhabitation}
%
A \emph{store} $\store$ is a mapping from
(free) variables to values.
%
We write $\apply{\store}{\expr}$ for the term
obtained by substituting free variables $x$
in $\expr$ with the corresponding values in
the store $\store$.
%
We write $\satisfies{\Renv}{\store}{\expr}$
if $\trans{\Renv}{\apply{\store}{\expr}}{\ttrue}$.
%
We say $\Penv$ is \emph{inhabited}
if for some $\store$ we have
$\satisfies{\Renv}{\store}{\pred}$
for each $\pred \in \Penv$.


\mypara{Divergence}
%
We say that a closed term $\expr$
\emph{diverges under} $\Renv$ if
there is no $\val$ such that
$\trans{\Renv}{\expr}{\val}$ .

\begin{lemma}[\textbf{Divergence}] \label{lem:div} % TRIVIAL
If $\ \forall i \in \nats$ we have
$\trans{\Renv}{\expr_i}{\ctxapp{\ctx}{\expr_{i+1}}}$
then $\expr_0$ diverges under $\Renv$.
\end{lemma}

\NV{What is the context?}

\mypara{Application Subterm}
We write $\issubterm{\expr}{\expr'}$
if $\expr$ is a \emph{syntactic subterm}
of $\expr'$ or if 
$\expr$ is a \textbf{fully applied} 
reflected function or data constructor application with arity $n$
($\expr \equiv f (e1, \dots, e_n)$)
and $\texttt{app}\ (\dots (\texttt{app}\ f\ e_1) \dots ) e_n$
is a \emph{syntactic subterm}
of $\expr'$.

\mypara{Subterm Evaluation}
%
The following lemma states that
if a store satisfies a reflected
function's guard, then the evaluation
of the corresponding function body
requires evaluating \emph{every subterm}
inside the body.

\begin{lemma}\label{lem:subterm-eval}
Let $\rlookup{\Renv}{\rfun} \equiv \rdef{x}{\pred}{\body}$.
%
If $\satisfies{\Renv}{\store}{\pred_i}$ and $\issubterm{\fapp{g}{e}}{\body_i}$
then
$\trans{\Renv}{\apply{\store}{\body_i}}
              {\ctxapp{\ctx}{\rawapp{g}{\apply{\store}{\params{\expr}}}}}$.
\end{lemma}

\NV{How is evaluation and context defined?}
\mypara{Totality Assumption}
%
The following assumption states that
if a store satisfies a reflected function's
guard, then the corresponding body reduces
to exactly one value.
%
Both of these are checked via
refinement typing \cite{Vazou14}.
%
\RJ{format below as "assumption" to avoid confusion}

\begin{lemma}[\textbf{Total}]\label{asm:total}
Let $\rlookup{\Renv}{\rfun} \equiv \rdef{x}{\pred}{\body}$.
%
$\rfun$ is \emph{total} if forall $\store$
%
\begin{enumerate}
\item If $\satisfies{\Renv}{\store}{\pred_i}$ then $\trans{\Renv}{\apply{\store}{\body_i}}{\val}$.
\item If $\satisfies{\Renv}{\store}{\pred_i}$, $\satisfies{\Renv}{\store}{\pred_j}$ then $i=j$.
\end{enumerate}
%
$\Renv$ is \emph{total} if every $\rfun \in \Renv$ is total.
%
\end{lemma}





\NV{The only interesting thing here is the transparency requirement. 
Yet, not clearly stated why you need it}

So far, we have shown that our proof search
procedure $\pbesym$ is both sound and complete.
%
Both of these are easy to achieve
simply by \emph{enumerating} all possible
instances and repeatedly querying the SMT
solver.
%
Such an approach --- lets call it
``monkeys-with-typewriters'' --- is rather
impractical: it may never terminate.
%
Fortunately, next, we show that in addition
to being sound and complete, our proof search
procedure always terminates.
%
The key insight is that if $\pbesym$ \emph{does not terminate}
then it is because
%
(1)~the set of depth $n$ unfoldings form a strictly
    increasing infinite chain,
    $\binst{\Renv}{\Penv}{0} \subset \binst{\Renv}{\Penv}{1} \ldots$,
    which means that
(2)~there is an infinite sequence of application terms
    $\stepsmany{\fapp{\rfun_0}{\expr_0}}{\stepsmany{\fapp{\rfun_1}{\expr_1}}{\ldots}}$
    which means that
(3)~the reflected function $\rfun_0$ is \emph{not total}, contradicting the
    totality assumption (\ref{asm:total}).
%
Thus, \pbesym must terminate!
%
Next, we formalize the above intuition and
prove the termination Theorem~\ref{thm:termination}.

\mypara{Symbolic Step}
%
A pair
$\symstep{\fapp{\rfun}{\expr}}{\fapp{\rfun'}{\expr'}}$
is a
$\Renv,\Penv$-\emph{symbolic step} (abbrev. step)
if
%
\begin{itemize}
  \item $\rlookup{\Renv}{\rfun} \equiv \rdef{x}{\pred}{\body}$,
  \item $\smtIsValid{\Penv \wedge Q}{\pred_i}$ for some $\Renv$-instance $Q$,
  \item $\issubterm{\fapp{\rfun'}{\expr'}}{\SUBSTS{\body_i}{x}{\expr}}$
\end{itemize}

\mypara{Steps and Reductions}
%
Next, using Lemmas~\ref{lem:subterm-eval},~\ref{lem:smt-approx}
and the definition of symbolic steps, we show that every
symbolic step under an inhabited $\Penv$ corresponds
to a \emph{sequence} of steps in the concrete semantics:

\begin{lemma}[\textbf{Step-Reductions}]\label{lem:step-reductions}
If   $\symstep{\fapp{\rfun}{\expr}}{\fapp{\rfun'}{\expr'}}$
     is a symbolic step under $\Renv, \Penv$
     and
     $\satisfies{\Renv}{\store}{\Penv}$
then $\trans{\Renv}
            {\fapp{\rfun}{\apply{\store}{\expr}}}
            {\ctxapp{\ctx}{\fapp{\rfun'}{\apply{\store}{\expr'}}}}$
for some context $\ctx$.
\end{lemma}

\mypara{Steps and Values}
%
The following lemma states that if
$\symstep{\fapp{\rfun}{y}}{\expr'}$
is a symbolic step under a $\Penv$
inhabited by $\store$ then
$\apply{\store}{\fapp{\rfun}{y}}$
reduces to a value.

\begin{lemma}[\textbf{Step-Value}]\label{lem:step-value}
If $\Renv$ is total,
   $\satisfies{\Renv}{\store}{\Penv}$, and
   $\symstep{\fapp{\rfun}{y}}{\expr'}$ is a $\Renv,\Penv$ step
then $\trans{\Renv}{\apply{\store}{\fapp{\rfun}{y}}}{\val}$.
\end{lemma}

The proof follows by observing that if $\Penv$
is inhabited by $\store$, and a particular step
is possible, then the guard corresponding to that
step must also be true under $\store$ (Lemma~\ref{lem:smt-approx})
and hence, by totality (\ref{asm:total}),
the function must reduce to a value under the given store.

\mypara{Symbolic Trace}
%
A sequence
%
$\fapp{\rfun_0}{\expr_0},\fapp{\rfun_1}{\expr_1},\fapp{\rfun_2}{\expr_2},\ldots$
%
is a $\Renv,\Penv$-\emph{symbolic trace} (abbrev. trace)
if $\params{\expr_0} \equiv \params{x_0}$ for some variables $x_0$,
and $\symstep{\fapp{\rfun_i}{\expr_i}}{\fapp{\rfun_{i+1}}{\expr_{i+1}}}$
is a $\Renv,\Penv$-step, for each $i$.

\begin{lemma}[\textbf{Finite-Trace}]\label{lem:finite-trace}
If $\Renv$ is total, $\satisfies{\Renv}{\store}{\Penv}$
then every $\Renv,\Penv$-trace is finite.
\end{lemma}

We prove the above lemma by contradiction.
Specifically, we show that an infinite
$\Renv,\Penv$-trace for a $\Penv$
inhabited by $\store$ yields an
infinite \emph{concrete trace}, which
contradicts totality, and hence, all the
$\Renv, \Penv$ symbolic traces
must be finite.

\mypara{Unfolding Chains and Traces}
%
If the unfolding $\Renv, \Penv$ yields
an infinite chain, then $\Renv, \Penv$
has an infinite symbolic trace:

\begin{lemma} [\textbf{Ascending Chains}] \label{lem:chain}
Let $\Penv_i \doteq \binst{\Renv}{\Penv}{i}$.
If there exists an (infinite) ascending chain
  $$\Penv_0 \subset \ldots \subset \Penv_n \ldots$$
then there exists an (infinite) symbolic trace
  $$\fapp{\rfun_0}{\expr_0},\ldots,\fapp{\rfun_n}{\expr_n},\ldots$$
\end{lemma}

To prove the above,
we construct a trace
by selecting at level $i$
(\ie in $\Penv_i$),
an application term
$\fapp{\rfun_i}{\expr_i}$
that was created by
unfolding an application
term at level $i-1$
(\ie in $\Penv_{i-1}$).


\mypara{Transparency}
%
We say $\Penv$ is \emph{inconsistent}
if $\smtIsValid{\Penv}{\tfalse}$.
%
We say $\Penv$ is \emph{transparent}
if it is either inhabited or inconsistent.
%
As an example of a \emph{non-transparent}
$\Penv_0$ consider the predicate
$\kw{lenA xs} = 1 + \kw{lenB xs}$,
where \kw{lenA} and \kw{lenB} are both
identical definitions of the list length
function.
%
Clearly there is no $\store$ that causes
the above predicate to evaluate to $\ttrue$.
%
At the same time, the SMT solver cannot
(using the decidable, quantifier-free theories)
prove a contradiction as that requires
induction over \kw{xs}.


\mypara{\pbesym Terminates}
%
Finally, we prove that as long the hypotheses
are \emph{transparent} the proof search
procedure $\pbesym$ terminates.

\begin{theorem}[\textbf{Termination}] \label{thm:termination}
  If $\Penv$ is transparent
  then $\pbe{\Renv}{\Penv}{\pred}$ terminates.
\end{theorem}

The proof has two cases.
%
If $\Penv$ is \emph{inconsistent} then trivially $\smtIsValid{\Penv}{\pred}$
and so $\pbe{\Renv}{\Penv}{\pred}$ terminates immediately.
%
If $\Penv$ is inhabited and yet $\pbesym$ loops forever,
there must be an infinite strictly ascending chain of
unfoldings $\Penv_i$, and hence, by lemma~\ref{lem:chain}
and infinite symbolic trace which contradicts
lemma~\ref{lem:finite-trace}.
%
Thus, if $\Penv$ is inhabited, $\pbesym$ must terminate.

%
The transparency requirement is not onerous.
In practice, the environments are either always
inconsistent \ie we have provably ``dead code'',
or it is easy to find concrete values that inhabit
the hypothesis by random~\cite{Claessen00}
or SMT-guided generation~\cite{Seidel15}.
